\uuid{3qGX}
\titre{Limite}
\auteur{bodin ; hanani ; mzari}
\organisation{exo7}
\licence{CC BY-NC-SA 4.0}
\datecreate{2020-01-20}
\niveau{L1}
\module{Analyse}
\chapitre{Suite}
\sousChapitre{Autre}
\difficulte{1}
\theme{Suites, limite}
\exo7qcmid{lille-1.suites.004}

\contenu{
\begin{question}
\qclassification{121.00}
\qtags{motcle=limite}
Soit $\displaystyle u_n=\sqrt{n+1}-\sqrt{n}$ et $\displaystyle v_n=\cos\left(\frac{n^2+1}{n^2-1}\pi\right)$. Quelles sont les bonnes réponses ?
\begin{answers}
    \bad{$\displaystyle \lim _{n\to +\infty}u_n=1$ et $\displaystyle \lim _{n\to +\infty}v_n=-1$ }
    \good{$\displaystyle \lim _{n\to +\infty}u_n=0$ et $\displaystyle \lim _{n\to +\infty}v_n=-1$}
    \bad{$\displaystyle \lim _{n\to +\infty}u_n=1$ et $\displaystyle \lim _{n\to +\infty}v_n=1$}
    \bad{$\displaystyle \lim _{n\to +\infty}u_n=0$ et $\displaystyle \lim _{n\to +\infty}v_n$ n'existe pas}
\end{answers}
\begin{explanations}
D'abord, $\displaystyle u_n=\frac{\left(\sqrt{n+1}-\sqrt{n}\right)\left(\sqrt{n+1}+\sqrt{n}\right)}{\sqrt{n+1}+\sqrt{n}}=\frac{1}{\sqrt{n+1}+\sqrt{n}}$ et donc $\displaystyle \lim _{n\to +\infty}u_n=0$ car $\displaystyle \lim _{n\to +\infty}\left(\sqrt{n+1}+\sqrt{n}\right)=+\infty$. Par ailleurs,
$$\displaystyle \frac{n^2+1}{n^2-1}\pi=\frac{n^2\left(1+\frac{1}{n^2}\right)}{n^2\left(1-\frac{1}{n^2}\right)}\pi=\frac{1+\frac{1}{n^2}}{1-\frac{1}{n^2}}\pi.$$
Donc $\displaystyle \lim _{n\to +\infty}\frac{n^2+1}{n^2-1}\pi=\pi$, et par suite, $\displaystyle \lim _{n\to +\infty}v_n=\cos (\pi)=-1$.
\end{explanations}
\end{question}
}
